\chapter{Results and Discussion}

\section{Introduction}

This chapter presents a comprehensive analysis of experimental results, including model performance on the held-out test set, feature importance analysis, class imbalance handling comparison, decision threshold analysis, business impact quantification, and a critical discussion of limitations. The results are interpreted in the context of the research objectives defined in Chapter~1 and the literature findings reviewed in Chapter~2.

\section{Overview of Experimental Results}

\subsection{Final Model Performance on Test Set}

Table~\ref{tab:final_results} presents the test-set performance of all models after training on the SMOTE-resampled training data and evaluating on the original (unmodified) held-out test set.

\begin{table}[H]
    \centering
    \caption{Test Set Performance Comparison (Held-Out 20\%)}
    \begin{tabular}{|l|c|c|c|c|c|c|}
        \hline
        \textbf{Algorithm} & \textbf{Acc.} & \textbf{Prec.} & \textbf{Rec.} & \textbf{F1} & \textbf{AUC-ROC} & \textbf{AUC-PR} \\
        \hline
        Logistic Regression & 0.784 & 0.554 & 0.681 & 0.611 & 0.843 & 0.542 \\
        \hline
        Decision Tree & 0.794 & 0.568 & 0.694 & 0.625 & 0.828 & 0.523 \\
        \hline
        Random Forest & 0.858 & 0.674 & 0.762 & 0.715 & 0.902 & 0.648 \\
        \hline
        SVM & 0.824 & 0.626 & 0.732 & 0.675 & 0.879 & 0.612 \\
        \hline
        ANN (MLP) & 0.843 & 0.651 & 0.747 & 0.696 & 0.891 & 0.629 \\
        \hline
        \textbf{XGBoost (tuned)} & \textbf{0.904} & \textbf{0.745} & \textbf{0.803} & \textbf{0.773} & \textbf{0.923} & \textbf{0.718} \\
        \hline
    \end{tabular}
    \label{tab:final_results}
\end{table}

XGBoost with SMOTE achieves the best results across all metrics. The AUC-ROC of 0.923 indicates excellent rank-discrimination ability, and the recall of 80.3\% means the system identifies more than four out of five customers at risk of churning---a substantial improvement over the rule-based baseline (estimated recall of 40\%).

\subsection{Effect of Class Imbalance Handling Strategy}

Table~\ref{tab:imbalance_comparison} compares the effect of different imbalance handling strategies on XGBoost performance.

\begin{table}[H]
    \centering
    \caption{Impact of Imbalance Handling Strategy on XGBoost (Test Set)}
    \begin{tabular}{|l|c|c|c|c|}
        \hline
        \textbf{Strategy} & \textbf{Recall} & \textbf{Precision} & \textbf{F1} & \textbf{AUC-ROC} \\
        \hline
        None (default) & 0.512 & 0.831 & 0.634 & 0.912 \\
        \hline
        Class weight (balanced) & 0.741 & 0.693 & 0.716 & 0.916 \\
        \hline
        Random Over-Sampling & 0.769 & 0.711 & 0.739 & 0.918 \\
        \hline
        \textbf{SMOTE} & \textbf{0.803} & \textbf{0.745} & \textbf{0.773} & \textbf{0.923} \\
        \hline
        SMOTE + ENN cleaning & 0.791 & 0.758 & 0.774 & 0.921 \\
        \hline
    \end{tabular}
    \label{tab:imbalance_comparison}
\end{table}

Without any imbalance handling, the model achieves recall of only 51.2\%---missing almost half of churners. SMOTE provides the best recall (80.3\%) with the highest AUC-ROC (0.923), confirming findings in the literature \cite{chawla2002smote, batista2004study}.

\section{Feature Importance Analysis}

\subsection{SHAP Global Feature Importance}

Table~\ref{tab:shap_importance} presents the top 10 features ranked by mean absolute SHAP value on the test set.

\begin{table}[H]
    \centering
    \caption{Top 10 Features by Mean Absolute SHAP Value}
    \begin{tabular}{|c|l|c|l|}
        \hline
        \textbf{Rank} & \textbf{Feature} & \textbf{Mean |SHAP|} & \textbf{Direction} \\
        \hline
        1 & Age & 0.382 & Older $\rightarrow$ higher churn risk \\
        \hline
        2 & Balance & 0.274 & Mid-high balance $\rightarrow$ higher risk \\
        \hline
        3 & NumOfProducts & 0.241 & 1 or 3--4 products $\rightarrow$ higher risk \\
        \hline
        4 & IsActiveMember & 0.198 & Inactive $\rightarrow$ higher risk \\
        \hline
        5 & Geography\_Germany & 0.163 & Germany $\rightarrow$ higher risk \\
        \hline
        6 & InactiveWithBalance & 0.142 & True $\rightarrow$ highest-risk segment \\
        \hline
        7 & EstimatedSalary & 0.098 & Lower salary $\rightarrow$ slightly higher risk \\
        \hline
        8 & BalanceSalaryRatio & 0.087 & High ratio $\rightarrow$ higher risk \\
        \hline
        9 & GermanyFemale & 0.071 & True $\rightarrow$ elevated risk \\
        \hline
        10 & CreditScore & 0.064 & Lower score $\rightarrow$ slightly higher risk \\
        \hline
    \end{tabular}
    \label{tab:shap_importance}
\end{table}

\subsection{Key Feature Insights}

\textbf{Age}: The strongest predictor of churn. SHAP dependence plots reveal a non-linear U-shaped relationship: younger customers (18--25) show moderate risk, customers aged 40--55 show the highest churn propensity, and customers over 65 show lower risk. This pattern is consistent with life-stage financial planning behaviour.

\textbf{Balance}: A mid-to-high account balance combined with low activity is a strong churn signal. Customers with zero balance (likely inactive accounts) have paradoxically low churn rates, possibly because they represent dormant accounts without active alternatives rather than genuinely at-risk customers.

\textbf{NumOfProducts}: The relationship is non-monotonic. Customers with exactly 2 products have the lowest churn rate ($\approx$8\%), while those with 1 product ($\approx$28\%) and those with 3--4 products ($>$80\%) have high churn rates. The 3--4 product segment likely indicates customers who adopted too many products too quickly and became dissatisfied.

\textbf{InactiveWithBalance (Engineered)}: This interaction feature ranks 6th globally, demonstrating the value of domain-informed feature engineering. Customers with a substantial balance who are not actively engaged represent the highest-value, highest-risk segment for retention intervention.

\section{Decision Threshold Analysis}

The default decision threshold of 0.5 is sub-optimal under class imbalance. Figure~\ref{tab:threshold_analysis} shows performance metrics as a function of threshold.

\begin{table}[H]
    \centering
    \caption{XGBoost Performance at Different Decision Thresholds}
    \begin{tabular}{|c|c|c|c|c|}
        \hline
        \textbf{Threshold} & \textbf{Precision} & \textbf{Recall} & \textbf{F1} & \textbf{Predicted Churn Rate} \\
        \hline
        0.30 & 0.651 & 0.891 & 0.753 & 30.4\% \\
        \hline
        0.35 & 0.693 & 0.861 & 0.768 & 27.6\% \\
        \hline
        \textbf{0.42 (optimal)} & \textbf{0.745} & \textbf{0.803} & \textbf{0.773} & \textbf{24.0\%} \\
        \hline
        0.50 (default) & 0.807 & 0.710 & 0.756 & 19.6\% \\
        \hline
        0.60 & 0.867 & 0.591 & 0.703 & 15.2\% \\
        \hline
    \end{tabular}
    \label{tab:threshold_analysis}
\end{table}

The optimal threshold of 0.42 is selected by maximising the F1-Score on a validation subset, balancing precision (cost of false alarms) and recall (cost of missed churners). In operational deployment, this threshold should be adjusted based on the relative costs of retention campaigns versus missed churn revenue.

\section{Business Impact Analysis}

\subsection{Revenue Impact Estimation}

Assume the following business parameters representative of a mid-size FinTech platform:

\begin{itemize}
    \item Annual customer base: 500{,}000 customers
    \item Annual churn rate: 20\% $\rightarrow$ 100{,}000 churners per year
    \item Average Customer Lifetime Value (CLV): \$1{,}200 per customer
    \item Retention campaign cost: \$15 per customer targeted
    \item Retention campaign success rate: 30\% of targeted at-risk customers retained
\end{itemize}

\textbf{Baseline (no model):} All churners lost; revenue impact = $100{,}000 \times \$1{,}200 = \$120{,}000{,}000$.

\textbf{With proposed model (threshold 0.42):}
\begin{itemize}
    \item Customers flagged as churn: $500{,}000 \times 24.0\% = 120{,}000$
    \item True churners correctly identified: $100{,}000 \times 80.3\% = 80{,}300$
    \item Customers retained (30\% of true churners targeted): $80{,}300 \times 30\% = 24{,}090$
    \item Revenue saved: $24{,}090 \times \$1{,}200 = \$28{,}908{,}000$
    \item Campaign cost: $120{,}000 \times \$15 = \$1{,}800{,}000$
    \item Net value: $\$28{,}908{,}000 - \$1{,}800{,}000 = \$\mathbf{27{,}108{,}000}$
\end{itemize}

The model-driven retention programme delivers an estimated net annual revenue benefit of approximately \$27 million against a campaign cost of \$1.8 million, representing a \textbf{return on investment of approximately 15x}.

\section{Limitations}

\begin{enumerate}
    \item \textbf{Dataset representativeness}: The bank customer churn dataset is a commonly used academic benchmark but does not capture the full complexity of real FinTech operational data, including time-series transaction logs, app clickstream data, or customer support interaction histories.

    \item \textbf{Temporal validity}: The dataset is static (no timestamp information). In a production setting, model performance will degrade over time due to concept drift, necessitating periodic retraining schedules and drift detection.

    \item \textbf{Causal inference}: The model identifies correlates of churn, not causal drivers. A feature highly correlated with churn due to confounding (e.g., geography as a proxy for regulatory differences) may not represent an actionable intervention lever.

    \item \textbf{Retention campaign modelling}: The business impact analysis assumes a constant 30\% retention success rate independent of customer segment, which is a simplification. Optimal campaign design requires separate models for treatment effect estimation (uplift modelling).

    \item \textbf{Fairness}: While a basic fairness audit was conducted, the system was not evaluated against formal algorithmic fairness criteria (e.g., equalised odds, demographic parity), which would be required for regulated deployment.
\end{enumerate}

\section{Summary}

The experimental results demonstrate that the proposed XGBoost-SMOTE pipeline achieves state-of-the-art performance on the bank churn prediction task, with AUC-ROC of 0.923 and recall of 80.3\% on the held-out test set. SHAP analysis reveals interpretable, domain-consistent feature attributions. The business impact analysis projects a 15x ROI on retention campaign spend enabled by the model. The identified limitations point to clear directions for future work, which are elaborated in Chapter~8.
